% !TeX root = ../main.tex

\chapter{Convergence}

% =================================================

We now investigate

\[
F(x)
=
\sum_{n=1}^{\infty}
c_n b^{x^n}.
\]

The behaviour depends strongly on both \(x\) and \(b\).

% -------------------------------------------------
\subsection{The Region \texorpdfstring{$|x|<1$}{|x|<1}}
% -------------------------------------------------

If

\[
|x|<1,
\]

then

\[
x^n\longrightarrow0.
\]

Consequently,

\[
b^{x^n}\longrightarrow1.
\]

Thus

\[
c_nb^{x^n}
\sim c_n.
\]

In particular, if

\[
c_n=1,
\]

then

\[
b^{x^n}\longrightarrow1,
\]

so the terms do not approach zero and

\[
\sum_{n=1}^{\infty}b^{x^n}
\]

diverges.

More generally, in the region \(|x|<1\), convergence is controlled
primarily by the coefficient sequence \(c_n\).

% -------------------------------------------------
\subsection{The Point \texorpdfstring{$x=1$}{x=1}}
% -------------------------------------------------

At

\[
x=1,
\]

we have

\[
x^n=1.
\]

Therefore

\[
F(1)
=
\sum_{n=1}^{\infty}c_nb.
\]

Hence

\[
\boxed{
F(1)=b\sum_{n=1}^{\infty}c_n.
}
\]

Thus convergence at \(x=1\) is exactly the same as convergence of

\[
\sum c_n.
\]

% -------------------------------------------------
\subsection{The Point \texorpdfstring{$x=0$}{x=0}}
% -------------------------------------------------

At

\[
x=0,
\]

we have

\[
0^n=0.
\]

Thus

\[
b^{0^n}=1
\]

and

\[
\boxed{
F(0)
=
\sum_{n=1}^{\infty}c_n.
}
\]

Again, convergence is determined directly by the coefficients.

% -------------------------------------------------
\subsection{The Region \texorpdfstring{$x>1$}{x>1} with \texorpdfstring{$b>1$}{b>1}}
% -------------------------------------------------

Suppose

\[
x>1
\]

and

\[
b>1.
\]

Then

\[
x^n\longrightarrow+\infty.
\]

Therefore

\[
b^{x^n}\longrightarrow+\infty.
\]

If \(c_n=1\), the terms themselves grow without bound.

For convergence we require at least

\[
\boxed{
c_nb^{x^n}\longrightarrow0.
}
\]

Thus the coefficients must decay extremely rapidly.

% -------------------------------------------------
\subsection{The Region \texorpdfstring{$x>1$}{x>1} with \texorpdfstring{$0<b<1$}{0<b<1}}
% -------------------------------------------------

Now suppose

\[
0<b<1.
\]

Then

\[
b^{x^n}\longrightarrow0
\]

extremely rapidly.

Consider the ratio

\[
\frac{b^{x^{n+1}}}{b^{x^n}}
=
b^{x^{n+1}-x^n}.
\]

Since

\[
x^{n+1}-x^n
=
x^n(x-1),
\]

we obtain

\[
\boxed{
\frac{b^{x^{n+1}}}{b^{x^n}}
=
b^{x^n(x-1)}.
}
\]

Because

\[
x^n(x-1)\to\infty
\]

and \(0<b<1\),

\[
b^{x^n(x-1)}
\longrightarrow0.
\]

Therefore

\[
\sum_{n=1}^{\infty}b^{x^n}
\]

converges for every \(x>1\).

The decay is faster than ordinary geometric decay.

% -------------------------------------------------
\subsection{Negative Values with \texorpdfstring{$x<-1$}{x<-1}}
% -------------------------------------------------

Negative values introduce a new phenomenon.

Let

\[
x=-r,
\qquad
r>1.
\]

Then

\[
x^n=(-1)^nr^n.
\]

Therefore, for even \(n\),

\[
x^n=r^n,
\]

while for odd \(n\),

\[
x^n=-r^n.
\]

Hence

\[
b^{x^n}
=
\begin{cases}
b^{r^n}, & n\text{ even},\\[4pt]
b^{-r^n}, & n\text{ odd}.
\end{cases}
\]

If

\[
b>1,
\]

the even-indexed blocks explode while the odd-indexed blocks decay.

If

\[
0<b<1,
\]

the opposite happens.

Thus one parity tends to zero while the other tends to infinity.

This creates a strong asymmetry that does not occur for \(x>0\).

% =================================================
\section{A Simple Example}
% =================================================

Consider

\[
b=\frac12
\]

and

\[
x=2.
\]

Then

\[
b^{x^n}
=
\left(\frac12\right)^{2^n}.
\]

The first terms are

\[
\frac{1}{2^2},
\qquad
\frac{1}{2^4},
\qquad
\frac{1}{2^8},
\qquad
\frac{1}{2^{16}},
\qquad
\frac{1}{2^{32}},
\ldots
\]

or

\[
\frac14,
\qquad
\frac1{16},
\qquad
\frac1{256},
\qquad
\frac1{65536},
\ldots.
\]

The exponent itself grows exponentially:

\[
2^n,
\]

and this exponential quantity appears inside another exponential.

This explains the extremely rapid decay.

% =================================================
